\section{An exact toroidal elliptic bridge to the candidate cusp}
\label{sec:teb-main}

This section asks for the smallest typed morphism between the rank-three
toroidal construction and the reconstructed Fabel--Alpöge threefold.  It also
tests whether that morphism is supplied by the nine-dimensional Hermitian Gram
cone.  The conclusions are distinct:
\begin{enumerate}
\item a quotient of integral lattices extends to a morphism of toric cusp
families and commutes with monodromy and the cusp filling maps;
\item the quotient has an explicit unimodular splitting as the candidate cusp
plus one elliptic cusp;
\item the order-three and order-four affine peripheral actions admit exact rank-six
lattice lifts;
\item none of these maps is induced by the determinant-rank stratification or
the rank-one \(\mathbb{CP}^{2}\) normal cones of the Gram construction.
\end{enumerate}
The last assertion is an obstruction to a structure-preserving
identification.  It is not the false assertion that no continuous map between
the underlying spheres exists.

\subsection{The candidate lattice and filling data}
\label{subsec:teb-candidate-data}

Retain the candidate map
\[
 f:X\longrightarrow B\simeq\mathbb P^{1},
 \qquad B^{\circ}=B\setminus\{p_{0},p_{1},p_{2}\}.
\]
Its smooth fibre has marked first homology
\[
 \Lambda_F=
 \mathbb Z\langle\widehat\gamma,\widehat u,
                       \widehat w,\widehat\delta\rangle.
\]
Put
\[
 A_F=\mathbb Z\langle\widehat w,\widehat\delta\rangle,
 \qquad
 B_F=\mathbb Z\langle\widehat\gamma,\widehat u\rangle.
\]
Thus \(\Lambda_F=A_F\oplus B_F\), with the decomposed order
\((\widehat w,\widehat\delta;\widehat\gamma,\widehat u)\).  Define
\[
 B_0=\begin{pmatrix}0&1\\-1&0\end{pmatrix}:B_F\longrightarrow A_F.
 \tag{TEB.1}
\]
The cusp monodromy is
\[
 M_F(a,b)=(a+B_0b,b),
 \tag{TEB.2}
\]
or, in the manuscript order
\((\widehat\gamma,\widehat u,\widehat w,\widehat\delta)\),
\[
 [M_F]=
 \begin{pmatrix}
 1&0&0&0\\
 0&1&0&0\\
 0&1&1&0\\
 -1&0&0&1
 \end{pmatrix}.
 \tag{TEB.3}
\]
In particular,
\[
 (M_F-I)^2=0,
 \qquad
 \ker(M_F-I)=\operatorname{im}(M_F-I)=A_F.
 \tag{TEB.4}
\]

The toric cusp neighbourhood retracts onto its central fibre \(W\).  In the
marking above, the induced filling homomorphism is
\[
 c_F:H_1(T^4;\mathbb Z)=A_F\oplus B_F\longrightarrow H_1(W;\mathbb Z)=B_F,
 \qquad c_F(a,b)=b.
 \tag{TEB.5}
\]
Thus the two toric angular cycles are killed and the two deck cycles survive.
The exact homology of the central fibre is
\[
 H_*(W;\mathbb Z)=
 (\mathbb Z,\mathbb Z^2,\mathbb Z^4,\mathbb Z^2,\mathbb Z).
 \tag{TEB.6}
\]
At the other two special points, the finite monodromies are
\[
 A_1=
 \begin{pmatrix}
 1&0&0&0\\6&0&1&0\\-6&-1&-1&0\\-2&1&0&1
 \end{pmatrix},
 \quad
 A_2=
 \begin{pmatrix}
 1&0&0&0\\0&0&-1&0\\-6&1&0&0\\3&0&1&1
 \end{pmatrix},
 \tag{TEB.7}
\]
in the manuscript order.  They satisfy
\[
 A_1^3=I,\qquad A_2^4=I,\qquad A_1A_2M_F=I.
 \tag{TEB.8}
\]
The invariant primitive vectors
\[
 v_1=\widehat\gamma+2\widehat u-4\widehat w,
 \qquad
 v_2=-\widehat\gamma-3\widehat u+3\widehat w
 \tag{TEB.9}
\]
define free affine actions
\[
 x\longmapsto A_jx+v_j/m_j,
 \qquad (m_1,m_2)=(3,4).
 \tag{TEB.10}
\]

\subsection{The rank-three cusp}
\label{subsec:teb-rank-three}

Let
\[
 L_3=\mathbb Z\langle e_1,e_2,e_3\rangle.
\]
For \(m\in L_3\) and \(\pi=(\pi_1,\pi_2,\pi_3)\in S_3\), put
\[
 \sigma_{m,\pi}=\operatorname{conv}\{m,
 m+e_{\pi_1},m+e_{\pi_1}+e_{\pi_2},m+e_1+e_2+e_3\}.
 \tag{TEB.11}
\]
These tetrahedra give the periodic Kuhn triangulation \(\mathcal T_3\) of
\(L_{3,\mathbb R}\).  Coning all cells at height one gives a smooth toric
fourfold \(Y_3\), because every maximal cone is unimodular.  The height
character gives
\[
 t:Y_3\longrightarrow\mathbb C,
\]
and a maximal affine chart has the exact equation
\[
 t=z_0z_1z_2z_3.
 \tag{TEB.12}
\]

For the untwisted quotient, a noncentral fibre is
\[
 \mathbb C^3/(s\mathbb Z^3+\mathbb Z^3),
 \qquad t=e^{2\pi i s},
\]
and its monodromy, in angular--deck order, is
\[
 M_3=\begin{pmatrix}I_3&I_3\\0&I_3\end{pmatrix},
 \qquad M_3(a,b)=(a+b,b).
 \tag{TEB.13}
\]
The toric filling map is
\[
 c_3:H_1(T^6;\mathbb Z)=L_3\oplus\overline L_3
 \longrightarrow\overline L_3,
 \qquad c_3(a,b)=b.                                    \tag{TEB.14}
\]
Indeed, each angular generator bounds a toric disc, while a deck generator
closes only by the corresponding lattice transformation and survives in the
filling group.

\subsection{The exact \texorpdfstring{\(A_2\)}{A2} quotient of the Kuhn fan}
\label{subsec:teb-fan-map}

In \(A_F\) define
\[
 n_1=\widehat w,\qquad
 n_2=\widehat\delta-\widehat w,\qquad
 n_3=-\widehat\delta.
 \tag{TEB.15}
\]
They are primitive and satisfy \(n_1+n_2+n_3=0\).  Define
\[
 P:L_3\longrightarrow A_F,
 \qquad P(e_i)=n_i.
 \tag{TEB.16}
\]
In the bases \((e_1,e_2,e_3)\) and
\((\widehat w,\widehat\delta)\),
\[
 P=\begin{pmatrix}1&-1&0\\0&1&-1\end{pmatrix},
 \qquad \ker P=\mathbb Z(1,1,1).
 \tag{TEB.17}
\]
Define the deck-lattice quotient
\[
 Q=B_0^{-1}P=
 \begin{pmatrix}0&-1&1\\1&-1&0\end{pmatrix}:
 \overline L_3\longrightarrow B_F.
 \tag{TEB.18}
\]
Then
\[
 B_0Q=P,\qquad \ker Q=\mathbb Z(1,1,1).
 \tag{TEB.19}
\]

\begin{proposition}[The Kuhn-to-\(A_2\) fan morphism]
\label{prop:teb-fan-morphism}
The map \(P\) sends every tetrahedron of \(\mathcal T_3\) onto a triangle of
the candidate's \(A_2\) triangulation.  Consequently
\[
 \widehat P:L_3\oplus\mathbb Z\longrightarrow A_F\oplus\mathbb Z,
 \qquad \widehat P(x,r)=(Px,r),
 \tag{TEB.20}
\]
is a map of height-one fans and induces a toric morphism
\[
 F:Y_3\longrightarrow Y_2,
 \qquad t\circ F=t.
 \tag{TEB.21}
\]
\end{proposition}

\begin{proof}
The four vertices of \(\sigma_{m,\pi}\) map to
\[
 Pm,\quad Pm+n_{\pi_1},\quad
 Pm+n_{\pi_1}+n_{\pi_2}=Pm-n_{\pi_3},\quad Pm.
 \tag{TEB.22}
\]
Write \((n_1,n_2,n_3)=(e_1,e_2-e_1,-e_2)\) in the target basis.  For the six
permutations, the nonrepeated images are
\[
\begin{array}{c|c}
(1,2,3)&\{v,v+e_1,v+e_2\}\\
(2,1,3)&\{v,v-e_1+e_2,v+e_2\}\\
(1,3,2)&\{v,v+e_1,v+e_1-e_2\}\\
(3,1,2)&\{v,v-e_2,v+e_1-e_2\}\\
(2,3,1)&\{v,v-e_1+e_2,v-e_1\}\\
(3,2,1)&\{v,v-e_2,v-e_1\}.
\end{array}                                             \tag{TEB.23}
\]
These are precisely the six triangles incident to \(v=Pm\) in the periodic
triangulation by
\[
 \operatorname{conv}\{u,u+e_1,u+e_2\},
 \quad
 \operatorname{conv}\{u+e_1,u+e_2,u+e_1+e_2\}.
\]
Thus every source cone maps into a target cone.  The toric fan criterion gives
\(F\); the last coordinate in (TEB.20) is unchanged, so the height character
is unchanged.
\end{proof}

The kernel \(\mathbb Z(1,1,1)\) is invariant under permutations of the three
coordinates.  Hence the quotient lattice is natural up to the Weyl action on
the \(A_2\) target.  The ordered basis and orientation used in (TEB.1) choose
one representative of that Weyl orbit.

\subsection{Monodromy and filling compatibility}
\label{subsec:teb-mono-filling}

\begin{theorem}[Exact local lattice bridge]
\label{thm:teb-local-lattice-bridge}
The homomorphism
\[
 \Phi=P\oplus Q:L_3\oplus\overline L_3
 \longrightarrow A_F\oplus B_F,
 \qquad \Phi(a,b)=(Pa,Qb),
 \tag{TEB.24}
\]
is surjective and obeys
\[
 \Phi M_3=M_F\Phi,
 \qquad c_F\Phi=Qc_3.                                  \tag{TEB.25}
\]
Thus the diagram
\[
\begin{CD}
H_1(T^6;\mathbb Z) @>{\Phi}>> H_1(T^4;\mathbb Z)\\
@V{c_3}VV @VV{c_F}V\\
\overline L_3 @>{Q}>> B_F
\end{CD}
\tag{TEB.26}
\]
commutes.
\end{theorem}

\begin{proof}
Both \(P\) and \(Q\) are surjective.  Using (TEB.13),
\[
 \Phi M_3(a,b)=(P(a+b),Qb).
\]
Using (TEB.2) and \(B_0Q=P\),
\[
 M_F\Phi(a,b)=(Pa+B_0Qb,Qb)=(Pa+Pb,Qb).
\]
This proves monodromy equivariance.  The filling identity is the literal
calculation
\[
 c_F\Phi(a,b)=Qb=Qc_3(a,b).
\]
\end{proof}

\subsection{The actual period correction and the cusp-family morphism}
\label{subsec:teb-period-lift}

The candidate period matrix at the cusp has the exact form
\[
 Z_2(t)=sB_0+C(t),
 \qquad t=e^{2\pi i s},
 \tag{TEB.27}
\]
where \(C(t)\) is holomorphic and bounded at \(t=0\).  Fix the integral right
inverse
\[
 S=\begin{pmatrix}1&1\\0&1\\0&0\end{pmatrix},
 \qquad PS=I_2,
 \tag{TEB.28}
\]
and define
\[
 C_3(t)=S C(t)Q,
 \qquad Z_3(t)=sI_3+C_3(t).
 \tag{TEB.29}
\]

\begin{theorem}[Holomorphic toroidal elliptic bridge]
\label{thm:teb-holomorphic-cusp-bridge}
After shrinking the cusp disc, \(\operatorname{Im}Z_3(t)\) is invertible and
\[
 PZ_3(t)=Z_2(t)Q.                                      \tag{TEB.30}
\]
Consequently \(P:\mathbb C^3\to\mathbb C^2\) descends on every punctured
fibre to a surjective holomorphic homomorphism
\[
 \frac{\mathbb C^3}{Z_3(t)\mathbb Z^3+\mathbb Z^3}
 \longrightarrow
 \frac{\mathbb C^2}{Z_2(t)\mathbb Z^2+\mathbb Z^2},    \tag{TEB.31}
\]
with connected elliptic kernel and integral homology map \(\Phi\).

Define the source and target toric actions by
\[
 \Psi^{(3)}_\lambda(x,t)=
 \bigl(\exp(2\pi iC_3(t)\lambda)t^\lambda x,t\bigr),
 \qquad \lambda\in\mathbb Z^3,                         \tag{TEB.32}
\]
\[
 \Psi^{(2)}_\mu(y,t)=
 \bigl(\exp(2\pi iC(t)\mu)t^{B_0\mu}y,t\bigr),
 \qquad \mu\in\mathbb Z^2.                            \tag{TEB.33}
\]
Then
\[
 F\Psi^{(3)}_\lambda=\Psi^{(2)}_{Q\lambda}F.          \tag{TEB.34}
\]
The actions are free and properly discontinuous on sufficiently small toric
cusp neighbourhoods, and (TEB.21) descends to a holomorphic map
\[
 \overline F:N_3^C\longrightarrow N_0                 \tag{TEB.35}
\]
over the cusp disc.  Its smooth-fibre and filling maps are (TEB.24) and
(TEB.26).
\end{theorem}

\begin{proof}
Equations (TEB.19), (TEB.28), and (TEB.29) give
\[
 PZ_3=sP+PSCQ=sB_0Q+CQ=Z_2Q.
\]
Because \(\operatorname{Im}s\to+\infty\) at the cusp while \(C_3(t)\) is
bounded, the determinant of \(\operatorname{Im}Z_3(t)\) has leading term
\((\operatorname{Im}s)^3\); it is nonzero on a sufficiently small punctured
disc.  Hence (TEB.31) is a well-defined homomorphism of complex tori.

Both kernel lattices in (TEB.17) and (TEB.19) are generated by
\(k=(1,1,1)\), and \(C_3(t)k=0\).  The kernel is therefore
\[
 \mathbb C k/(\mathbb Z k+s\mathbb Z k),
\]
an elliptic curve.  Saturation of the primitive kernel proves connectedness.

On the dense torus, (TEB.30) gives
\[
\begin{aligned}
 F\Psi^{(3)}_\lambda(x,t)
 &=\bigl(\exp(2\pi iPC_3(t)\lambda)t^{P\lambda}F(x),t\bigr)\\
 &=\bigl(\exp(2\pi iC(t)Q\lambda)t^{B_0Q\lambda}F(x),t\bigr),
\end{aligned}
\]
which is (TEB.34).  Proposition~\ref{prop:teb-fan-morphism} extends the first
map over every toric stratum.

For completeness, freeness at a boundary point follows from orbit labels.  A
nonzero \(\lambda\) translates the bounded cell labelling its toric orbit to
that cell plus \(\lambda\); equality would force \(\lambda=0\).  On a smooth
fibre a fixed point would give
\(Z_3(t)\lambda\in\mathbb Z^3\), and taking imaginary parts gives
\(\operatorname{Im}Z_3(t)\lambda=0\), again forcing \(\lambda=0\).
Here is the required uniform estimate for proper discontinuity and
properness.  Put \(y=\operatorname{Im}s=-\log|t|/(2\pi)\), and, on a
punctured fibre, use
\[
 \operatorname{trop}_t(x)
 =\frac{1}{\log|t|}
   \bigl(\log|x_1|,\log|x_2|,\log|x_3|\bigr).
\]
The exact transformation law is
\[
 \operatorname{trop}_t(\Psi^{(3)}_\lambda x)
 =\operatorname{trop}_t(x)+A(t)\lambda,\qquad
 A(t)=I_3+\frac{\operatorname{Im}C_3(t)}{y}.           \tag{TEB.35a}
\]
Let
\[
 K=\sup_{|t|\leq\epsilon}
   \|\operatorname{Im}C_3(t)\|_{\mathrm{op}}<\infty.
\]
After shrinking \(\epsilon\), \(K/y\leq1/2\), and hence
\[
 \frac12\|\lambda\|
 \leq\|A(t)\lambda\|
 \leq\frac32\|\lambda\|                               \tag{TEB.35b}
\]
uniformly.  If a compact set meets its \(\lambda\)-translate in the dense
locus, its tropical coordinates lie in one bounded set \(C\), so
\(\|A(t)\lambda\|\leq\operatorname{diam}(C)\) and (TEB.35b) leaves only
finitely many integral \(\lambda\).  On the toric boundary the cell label is
translated by \(\lambda\), because the nonzero factor
\(\exp(2\pi iC_3(t)\lambda)\) does not change an orbit label; finitely many
cells again allow only finitely many \(\lambda\).

For properness, given a tropical point \(r\), choose
\(\lambda\in\mathbb Z^3\) by rounding the three coordinates of
\(A(t)^{-1}r\).  Then \(r-A(t)\lambda\) lies in
\(A(t)[0,1]^3\), a uniformly bounded set by (TEB.35b).  Over a smaller closed
disc, the closures of these representatives lie in the finite union of the
closed toric polydiscs belonging to the six Kuhn tetrahedra, exactly as in
the untwisted rank-three quotient.  That union is compact.  Thus the action
is properly discontinuous, the inverse image of every smaller closed disc
has compact quotient, and the quotient and descended map in (TEB.35) exist.
The homology assertions follow from
Theorem~\ref{thm:teb-local-lattice-bridge}.
\end{proof}

\subsection{A simultaneous unimodular splitting}
\label{subsec:teb-conjugacy}

\begin{proposition}[Candidate cusp plus one elliptic cusp]
\label{prop:teb-explicit-conjugacy}
Let
\[
 r=\begin{pmatrix}0&0&1\end{pmatrix},
\]
\[
 A=\begin{pmatrix}P\\r\end{pmatrix}
 =\begin{pmatrix}1&-1&0\\0&1&-1\\0&0&1\end{pmatrix},
 \qquad
 B=\begin{pmatrix}Q\\r\end{pmatrix}
 =\begin{pmatrix}0&-1&1\\1&-1&0\\0&0&1\end{pmatrix},
 \tag{TEB.36}
\]
and
\[
 D=B_0\oplus[1]
 =\begin{pmatrix}0&1&0\\-1&0&0\\0&0&1\end{pmatrix}.
 \tag{TEB.37}
\]
Then
\[
 \det A=\det B=1,\qquad A=DB.                           \tag{TEB.38}
\]
For \(J=\operatorname{diag}(A,B)\in SL(6,\mathbb Z)\),
\[
 JM_3J^{-1}=
 \begin{pmatrix}I_3&D\\0&I_3\end{pmatrix}.            \tag{TEB.39}
\]
In the order \((A_F,k_A;B_F,k_B)\), this is the literal matrix
\[
\begin{pmatrix}
1&0&0&0&1&0\\
0&1&0&-1&0&0\\
0&0&1&0&0&1\\
0&0&0&1&0&0\\
0&0&0&0&1&0\\
0&0&0&0&0&1
\end{pmatrix}.
\]
In the coordinates
\[
 (A_F,k_A;B_F,k_B)=(Pa,ra;Qb,rb),
\]
and after reordering them as
\[
 (A_F,B_F;k_A,k_B),
\]
the right side of (TEB.39) is
 \[
 M_F\oplus C_0,
 \qquad C_0=\begin{pmatrix}1&1\\0&1\end{pmatrix}.      \tag{TEB.40}
 \]
Explicitly, in the reordered basis it is
\[
\begin{pmatrix}
1&0&0&1&0&0\\
0&1&-1&0&0&0\\
0&0&1&0&0&0\\
0&0&0&1&0&0\\
0&0&0&0&1&1\\
0&0&0&0&0&1
\end{pmatrix}
=
\begin{pmatrix}M_F&0\\0&C_0\end{pmatrix},
\]
where \(M_F\) on the right is in the decomposed order
\((A_F;B_F)\).
Projection to the candidate block is exactly \(\Phi=P\oplus Q\).
\end{proposition}

\begin{proof}
The matrix \(A\) is upper triangular with unit diagonal.  Expanding \(B\)
along its last row gives
\[
 \det B=\det\begin{pmatrix}0&-1\\1&-1\end{pmatrix}=1.
\]
Direct multiplication gives \(DB=A\).  Hence
\[
 JM_3J^{-1}
 =\begin{pmatrix}A&0\\0&B\end{pmatrix}
  \begin{pmatrix}I&I\\0&I\end{pmatrix}
  \begin{pmatrix}A^{-1}&0\\0&B^{-1}\end{pmatrix}
 =\begin{pmatrix}I&AB^{-1}\\0&I\end{pmatrix}
 =\begin{pmatrix}I&D\\0&I\end{pmatrix}.
\]
The first two rows of \(A\) and \(B\) are \(P\) and \(Q\); their last row
records the primitive diagonal kernel coordinate.  Since \(D=B_0\oplus1\),
the asserted permutation separates the candidate and kernel blocks exactly.
\end{proof}

This proposition proves integrally, rather than only rationally, that the
rank-three cusp module is the candidate cusp module plus one standard rank-one
cusp.  It also identifies both kernel cycles with the primitive diagonal
vector \((1,1,1)\).

\subsection{The complete finite-order lattice lift}
\label{subsec:teb-finite-lift}

Define
\[
 C_1=\begin{pmatrix}-1&-1\\1&0\end{pmatrix},
 \qquad
 C_2=\begin{pmatrix}0&1\\-1&0\end{pmatrix},
 \qquad
 C_0=\begin{pmatrix}1&1\\0&1\end{pmatrix}.
 \tag{TEB.41}
\]
Direct multiplication gives
\[
 C_1^3=I,\qquad C_2^4=I,\qquad
 C_1C_2=C_0^{-1},\qquad
 (C_0-I)^2=0,\quad\operatorname{rank}(C_0-I)=1.
 \tag{TEB.42}
\]

\begin{proposition}[Rank-six lift of the \((3,4,\infty)\) peripheral actions]
\label{prop:teb-finite-monodromy-lift}
On \(\widetilde\Lambda=\Lambda_F\oplus\mathbb Z^2\), put
\[
 \widetilde A_1=A_1\oplus C_1,
 \qquad
 \widetilde A_2=A_2\oplus C_2,
 \qquad
 \widetilde M_0=M_F\oplus C_0.
 \tag{TEB.43}
\]
Then
\[
 \widetilde A_1^3=I,\qquad
 \widetilde A_2^4=I,\qquad
 \widetilde A_1\widetilde A_2\widetilde M_0=I,          \tag{TEB.44}
\]
and \(\widetilde M_0-I\) has square zero, rank three, and equal kernel and
image of rank three.  Projection
\[
 \operatorname{pr}:\widetilde\Lambda\longrightarrow\Lambda_F
\]
intertwines all three monodromies.

The affine maps
\[
 \widetilde g_j(x,y)=
 \bigl(A_jx+v_j/m_j,C_jy\bigr)
 \quad\text{on }T^4\times T^2                          \tag{TEB.45}
\]
 have orders \(m_1=3,m_2=4\) and are free.  Their quotient maps project to the
candidate finite-fibre quotient models.
\end{proposition}

\begin{proof}
Equations (TEB.8) and (TEB.42) prove (TEB.44) blockwise.  The nilpotent part
of \(\widetilde M_0\) is the direct sum of the candidate rank-two nilpotent
and the rank-one nilpotent \(C_0-I\), which proves all cusp assertions.
Block projection gives equivariance.

The vectors \((v_j,0)\) are invariant under \(\widetilde A_j\), so the
\(m_j\)-th iterate of (TEB.45) is translation by the integral vector
\((v_j,0)\), hence the identity on the torus.  If a nonidentity iterate had a
fixed point, its projection to \(T^4\) would be a fixed point of the
corresponding nonidentity iterate of (TEB.10).  The latter actions are free.
Thus (TEB.45) is free and has exact order \(m_j\).  Projection descends because
it is equivariant.
\end{proof}

Proposition~\ref{prop:teb-finite-monodromy-lift} proves that the orders
\(3,4\), affine freeness, and a rank-three cusp are mutually compatible at
the integral topological level.  It does not yet identify these block lifts
with one holomorphic period family on the full triangle-group cover.  That is
part of the remaining problem below.

\subsection{The six toric chambers and the diagonal Gram boundary}
\label{subsec:teb-weyl-chambers}

The numerical agreement between the six deepest toric strata and the six
\(A_2\) Weyl chambers has the following exact, but purely
fan-theoretic, form.  Let
\[
 \mathfrak a=
 \left\{\operatorname{diag}(\lambda_1,\lambda_2,\lambda_3):
 \lambda_1+\lambda_2+\lambda_3=0\right\}
 \subset\operatorname{Herm}_3^0(\mathbb C)
\]
and put
\[
 \nu_i=E_{ii}-\frac13I_3,\qquad
 \nu_1+\nu_2+\nu_3=0.                                  \tag{TEB.45a}
\]
Because \(n_1+n_2+n_3=0\), there is a unique real-linear isomorphism
\[
 \jmath:A_{F,\mathbb R}\xrightarrow{\ \sim\ }\mathfrak a,
 \qquad \jmath(n_i)=\nu_i.                             \tag{TEB.45b}
\]
For \(H\in\operatorname{Herm}_3^0(\mathbb C)\), let
\[
 \Phi_{\mathrm G}(H)=H-\lambda_{\min}(H)I_3.
\]
This is the \(SU(3)\)-equivariant homeomorphism onto the singular
positive-semidefinite Gram cone used in the preceding higher-rung analysis.

\begin{proposition}[The Kuhn strata are the marked \(A_2\) chambers]
\label{prop:teb-weyl-chambers}
For every \(\pi=(\pi_1,\pi_2,\pi_3)\in S_3\),
\[
 \jmath\!\left(P(\sigma_{0,\pi})\right)
 =
 \operatorname{conv}
 \{0,\nu_{\pi_1},-\nu_{\pi_3}\},                       \tag{TEB.45c}
\]
and the cone on this triangle is the closed Weyl chamber
\[
 \boxed{
 \mathbb R_{\geq0}\nu_{\pi_1}
 +\mathbb R_{\geq0}(-\nu_{\pi_3})
 =
 \{\lambda_{\pi_1}\geq\lambda_{\pi_2}
                  \geq\lambda_{\pi_3}\}.}              \tag{TEB.45d}
\]
Consequently the six tetrahedron orbits, hence the six quadruple-point
strata of the untwisted toroidal central fibre, are identified
\(S_3\)-equivariantly with the six marked diagonal Weyl chambers.

On this chamber the Gram-boundary homeomorphism is the explicit
piecewise-linear map
\[
 \Phi_{\mathrm G}(H)=H-\lambda_{\pi_3}(H)I_3.          \tag{TEB.45e}
\]
In particular,
\[
 \Phi_{\mathrm G}(\nu_{\pi_1})=E_{\pi_1\pi_1},
 \qquad
 \Phi_{\mathrm G}(-\nu_{\pi_3})=I_3-E_{\pi_3\pi_3}.    \tag{TEB.45f}
\]
\end{proposition}

\begin{proof}
Equation (TEB.22) with \(m=0\) gives the three distinct image vertices
\(0,n_{\pi_1},-n_{\pi_3}\), proving (TEB.45c).  For
\(a,b\geq0\), the diagonal entries of
\(a\nu_{\pi_1}-b\nu_{\pi_3}\) satisfy
\[
 \lambda_{\pi_1}-\lambda_{\pi_2}=a,\qquad
 \lambda_{\pi_2}-\lambda_{\pi_3}=b.
\]
Conversely these two nonnegative differences recover \(a,b\), since the
trace is zero.  This proves (TEB.45d).  Both the permutation action on the
Kuhn labels and the Weyl action on diagonal entries are the same action of
\(S_3\), so the correspondence is equivariant.  On the displayed chamber
the least eigenvalue is \(\lambda_{\pi_3}\), which proves (TEB.45e).
Substitution of the two extreme rays gives (TEB.45f).
\end{proof}

This is stronger than the bare equality \(6=|S_3|\), but it is not an
identification of singular spaces.  After quotienting the target fan by
\(A_F\), the six triangles incident to a marked vertex form two global
triangle orbits.  Moreover \(\Phi_{\mathrm G}\circ\jmath\) is a map on the
real cocharacter space, not a holomorphic map of toroidal total spaces.
It neither supplies the period correction nor changes the analytic
factorization comparison below.

\subsection{The exact Jordan determinant and Lorentz sections}
\label{subsec:teb-jordan-lorentz-sections}

The determinant--index identity in
Corollary~\ref{hrftd:cor:euler-shear-determinant} has a precise extension
inside the nine-dimensional Hermitian cone.  It is an algebraic extension of
the scalar invariant; it is not a map from the toroidal degeneration to the
Hermitian cone.

Write
\[
 \mathcal J_3=\operatorname{Herm}_3(\mathbb C),\qquad
 \mathcal J_3^{\geq0}=\{H\in\mathcal J_3:H\succeq0\},
 \qquad \Delta(H)=\det H.
\]
On the ambient vector spaces, the inclusion
\[
 \iota:\operatorname{Sym}_3(\mathbb R)\hookrightarrow\mathcal J_3,
 \qquad \iota(S)=S+i0                                      \tag{TEB.45g}
\]
is an injective unital homomorphism of real Jordan algebras for the product
\(A\circ B=(AB+BA)/2\).  Its restriction sends the positive-definite cone
into the positive-definite Hermitian cone.  The ambient formulation matters:
a positive cone is not itself a vector space or a Jordan algebra.  The map
also uses the retained marking of the integral lattice; it is not a
coordinate-free map from an unmarked congruence class of polarizations to a
fixed matrix space.

For \(a=(a_1,a_2,a_3)^{\mathsf T}\in\mathbb R^3\), fix the oriented
cross-product matrix
\[
 [a]_{\times}=
 \begin{pmatrix}
  0&-a_3&a_2\\
  a_3&0&-a_1\\
  -a_2&a_1&0
 \end{pmatrix}.                                           \tag{TEB.45h}
\]
Every \(H\in\mathcal J_3\) has one and only one decomposition
\[
 H=S+i[a]_{\times},\qquad
 S\in\operatorname{Sym}_3(\mathbb R),\quad a\in\mathbb R^3. \tag{TEB.45i}
\]

\begin{theorem}[Hermitian determinant splitting]
\label{thm:teb-hermitian-determinant-splitting}
For every real symmetric \(S\) and every \(a\in\mathbb R^3\),
\[
 \boxed{\det\bigl(S+i[a]_{\times}\bigr)
        =\det S-a^{\mathsf T}Sa.}                         \tag{TEB.45j}
\]
If \(S\) is positive definite, then
\[
 S+i[a]_{\times}\succeq0
 \quad\Longleftrightarrow\quad
 a^{\mathsf T}Sa\leq\det S,                              \tag{TEB.45k}
\]
and positive semidefiniteness and the weak scalar inequality may both be
replaced by strict positivity and a strict scalar inequality.
\end{theorem}

\begin{proof}
Direct expansion in the fixed coordinates (TEB.45h) gives
\[
 \det\bigl(S+i[a]_{\times}\bigr)=\det S-a^{\mathsf T}Sa.
\]
For the order statement let \(M=S^{-1/2}\), the positive-definite symmetric
inverse square root.  The identity
\[
 M[a]_{\times}M^{\mathsf T}
 =\det(M)[M^{-\mathsf T}a]_{\times}                       \tag{TEB.45l}
\]
follows by applying both sides to two vectors and evaluating against a third
with the fixed oriented volume form.  Hence
\[
 S^{-1/2}\bigl(S+i[a]_{\times}\bigr)S^{-1/2}
 =I_3+i[b]_{\times},
 \qquad b=\det(S)^{-1/2}S^{1/2}a,
\]
and
\[
 \lVert b\rVert^2=\frac{a^{\mathsf T}Sa}{\det S}.
\]
The three eigenvalues of \(I_3+i[b]_{\times}\) are
\(1-\lVert b\rVert,1,1+\lVert b\rVert\).  Congruence preserves inertia,
so (TEB.45k) and its strict version follow.
\end{proof}

For a fixed marked positive-definite symmetric monodromy form \(B\), define
the injective real-linear map
\[
 \ell_B:\mathbb R\oplus\mathbb R^3\longrightarrow\mathcal J_3,
 \qquad \ell_B(\lambda,a)=\lambda B+i[a]_{\times}.        \tag{TEB.45m}
\]
Theorem~\ref{thm:teb-hermitian-determinant-splitting}, applied first to
\(\lambda>0\) and then directly at \(\lambda=0\), gives
\[
 \det\ell_B(\lambda,a)
 =\lambda\bigl(\lambda^2\det B-a^{\mathsf T}Ba\bigr),    \tag{TEB.45n}
\]
and the exact inverse image
\[
 \boxed{
\ell_B^{-1}(\mathcal J_3^{\geq0})
=\{(\lambda,a):\lambda\geq0,
       \ a^{\mathsf T}Ba\leq\lambda^2\det B\}.}         \tag{TEB.45o}
\]
Indeed, if \(\lambda<0\), then
\(\operatorname{tr}\ell_B(\lambda,a)=\lambda\operatorname{tr}B<0\),
which is impossible for a positive-semidefinite matrix.
Thus every marked positive monodromy ray determines a literal future
\((3+1)\)-Lorentz cone with quadratic form
\[
 q_B(\lambda,a)=\lambda^2\det B-a^{\mathsf T}Ba.          \tag{TEB.45p}
\]
At \(B=I_3\) this is
\(\lambda^2-a_1^2-a_2^2-a_3^2\).  The equality
\[
 \chi(W_{3,I_3})=6=6q_{I_3}(1,0)                         \tag{TEB.45q}
\]
is therefore exact, while remaining an equality of scalar invariants rather
than a morphism of degenerations.

There is also an exact order-cone lift of the rank-two Lorentz cone from
Theorem~\ref{thm:hr3c-two-ladders}.  In its coordinates
\[
 H_2(t,x_1,x_2,x_3)=
 \begin{pmatrix}
  t+x_3&x_1-i x_2\\
  x_1+i x_2&t-x_3
 \end{pmatrix},
\]
put
\[
 \Theta(H_2(t,x_1,x_2,x_3))=
 \begin{pmatrix}
  t&-i x_3&-i x_2\\
  i x_3&t&-i x_1\\
  i x_2&i x_1&t
 \end{pmatrix}.                                         \tag{TEB.45r}
\]

\begin{proposition}[Order-cone lift and its complete rank behaviour]
\label{prop:teb-theta-order-lift}
The map \(\Theta\) is a real-linear bijection from
\(\operatorname{Herm}_2(\mathbb C)\) onto
\(\mathbb RI_3\oplus i\mathfrak{so}(3)\), and
\[
 H_2\succeq0\quad\Longleftrightarrow\quad
 \Theta(H_2)\succeq0.                                  \tag{TEB.45s}
\]
Moreover
\[
 \boxed{\det_3\Theta(H_2)
 =\frac{\operatorname{tr}H_2}{2}\det_2H_2.}             \tag{TEB.45t}
\]
On the positive-semidefinite cone, \(\Theta\) sends ranks
\(2,1,0\) respectively to ranks \(3,2,0\).
\end{proposition}

\begin{proof}
Put \(r=(x_1^2+x_2^2+x_3^2)^{1/2}\).  The eigenvalues of \(H_2\) are
\(t-r,t+r\), while those of \(\Theta(H_2)\) are
\(t-r,t,t+r\).  Both matrices are positive semidefinite exactly when
\(t\geq r\).  Multiplication of the displayed eigenvalues gives
\[
 \det_3\Theta(H_2)=t(t^2-r^2)
 =\frac{\operatorname{tr}H_2}{2}\det_2H_2.
\]
The same eigenvalues give all three rank assertions.  The displayed matrix
also makes real-linearity, injectivity, and the stated image immediate.
\end{proof}

If \(P\) is a rank-one projection in
\(\operatorname{Herm}_2(\mathbb C)\), then its Lorentz coordinates satisfy
\(t=r=1/2\), so \(\Theta(P)\) has eigenvalues \(0,1/2,1\) and is not an
idempotent.  Thus \(\Theta\) is not a Jordan-algebra homomorphism.  Globally,
the pullback of the cubic determinant hypersurface is
\[
 \Theta^{-1}\{\det_3=0\}
 =\{\operatorname{tr}H_2=0\}\cup\{\det_2H_2=0\};         \tag{TEB.45u}
\]
only on the positive cone does the first set add no points beyond the origin.
Neither \(\ell_B\) nor \(\Theta\) is a holomorphic toric map, a map of
integral period lattices, or a singular-fibre filling map.  For the actual
one-parameter toric degeneration the monodromy form is fixed at \(B=I_3\)
while the local parameter in \(z_0z_1z_2z_3=t\) tends to zero.  Consequently
the induced Jordan datum is the constant point \(I_3\), not a path to
\(\det B=0\).  Degeneration at fixed monodromy and rank loss of the monodromy
form are different operations.

Finally, the cyclic group \(C_3\subset S_3\) does have two orbits on the six
Kuhn labels.  Under \(P\), the even labels
\((123),(231),(312)\) give translates of the lower target triangle
\(\operatorname{conv}\{u,u+e_1,u+e_2\}\), while the odd labels
\((132),(213),(321)\) give translates of the upper target triangle
\(\operatorname{conv}\{u+e_1,u+e_2,u+e_1+e_2\}\).  This action is by
coordinate permutations and the exact bridge lowers lattice rank through
\(P\); it is not translation by an index-three superlattice.  No such
translation-lattice realization is constructed by the packet supplying
(TEB.45j)--(TEB.45t).  More decisively,
for any rank-three lattice \(\Lambda'\), a fundamental domain has normalized
volume \(3!=6\), so every \(\Lambda'\)-unimodular triangulation of that domain
has six maximal tetrahedra.  The canonical six-to-two construction is instead
the rank-lowering fan map \(P:L_3\to A_F\) in (TEB.16)--(TEB.23).  It must not
be replaced by an unproved semistable translation quotient.

\begin{auditrecord}[Completed Jordan--toric packet, 3 September 2026]
The source archive has SHA-256
\Source{8BA0BD7FD7BADEAAF2C44E2F2B4EFD42C25DA9119DA3180017FC488CA364E96F}.
Its supplied symbolic identities were replayed independently, and separate
Euler/stratification, morphism/type, fan/Weyl, and theorem-map audits were
performed.  Equations (TEB.45j)--(TEB.45t) are the surviving exact content.
The positive-cone Jordan typing, the normalization stratum attribution, and
the index-three compression were corrected rather than copied.  The audit
reports and exact replay script are retained under \Source{agent_logs/} and
\Source{agent_scratch/jordan_toric_symbolic_replay_20260903/} in the local
research record.
\end{auditrecord}

\subsection{The Gram determinant does not supply this bridge}
\label{subsec:teb-gram-obstructions}

Let
\[
 V=\operatorname{Mat}_3(\mathbb C),
 \qquad
 \mathcal G=V^+/SU(3),
 \qquad
 \mathcal D=\{\det=0\}^+/SU(3).
\]
Here \(V^+=V\cup\{\infty\}\) is the one-point compactification and
\(\{\det=0\}^+=\{\det=0\}\cup\{\infty\}\).
The exact Gram classification gives
\[
 (\mathcal G,\mathcal D)\cong(S^{10},S^8),
 \qquad
 \mathcal G\setminus\mathcal D\cong\mathbb R^9\times S^1.
 \tag{TEB.46}
\]
The real fixed form gives
\[
 \left(\operatorname{Mat}_3(\mathbb R)^+/SO(3),
       \{\det=0\}^+/SO(3)\right)
 \cong(S^6,S^5).                                       \tag{TEB.47}
\]

\begin{theorem}[No determinant-stratified identification]
\label{thm:teb-no-gram-identification}
There is no homeomorphism of pairs from the real Gram pair (TEB.47) to
\[
 \left(X,f^{-1}\{p_0,p_1,p_2\}\right).
 \tag{TEB.48}
\]
There is no homotopy equivalence from the complex Gram complement in
(TEB.46) to the candidate regular part \(f^{-1}(B^\circ)\).  The rank-one
trace-one Gram stratum \(\mathbb{CP}^2\) is not homotopy equivalent to any of
the three candidate special fibres.  Finally, no equivalence of function
germs identifies its normal determinant with any of the three candidate local
filling maps.
\end{theorem}

\begin{proof}
The determinant-zero set in (TEB.47) is a connected codimension-one
\(S^5\), and its complement has two components.  The set in (TEB.48) is the
disjoint union of three connected real four-dimensional fibres, hence has
codimension two and three components.  Its complement is connected because
it is a torus bundle over the connected pair of pants \(B^\circ\).  This
proves the first assertion.

The complex Gram complement has fundamental group \(\mathbb Z\).  The
homotopy sequence of the candidate regular torus bundle gives
\[
 1\longrightarrow\mathbb Z^4\longrightarrow
 \pi_1\bigl(f^{-1}(B^\circ)\bigr)
 \longrightarrow\pi_1(B^\circ)=F_2\longrightarrow1.    \tag{TEB.49}
\]
Its fundamental group surjects onto the nonabelian group \(F_2\), so it is
not \(\mathbb Z\).

The rank-one trace-one stratum has
\[
 H_*(\mathbb{CP}^2;\mathbb Z)=(\mathbb Z,0,\mathbb Z,0,\mathbb Z),
 \qquad \chi(\mathbb{CP}^2)=3.
\]
The cusp fibre has (TEB.6), so \(b_1=2,b_2=4,\chi=2\).  Each finite reduced
fibre is a free finite quotient of \(T^4\), hence has \(\chi=0\), and its
audited first homology has free rank two.  These invariants rule out all three
homotopy equivalences.

At a rank-one Gram matrix, Schur-complement coordinates put the normal
determinant in the exact form
\[
 uv-x^2-y^2,
 \tag{TEB.50}
\]
a nondegenerate quadratic form of signature \((1,3)\).  Its complexification
is the irreducible quadratic determinant
\[
 s_{11}s_{22}-s_{12}s_{21}.
 \tag{TEB.51}
\]
At a deepest candidate cusp point, the map is instead
\[
 (z_0,z_1,z_2)\longmapsto z_0z_1z_2,                   \tag{TEB.52}
\]
with three smooth complex branches.  The source dimensions, multiplicities,
and number of irreducible branches in (TEB.50)--(TEB.52) disagree.  The two
finite transverse maps have multiplicities three and four, while the Gram
determinant has multiplicity two along the rank-one stratum and multiplicity
three at the zero matrix.  It has no order-four stratum.  Orders of vanishing
and analytic factorization are invariant under equivalence of function germs,
which proves the final assertion.
\end{proof}

There are two further typed obstructions.  First, the determinant phase map
\[
 \arg\det:\mathcal G\setminus\mathcal D\longrightarrow S^1
\]
has trivial page monodromy after the \(SU(3)\) quotient.  Before quotienting,
its page is diffeomorphic to
\(\mathbb R_{>0}\times SL(3,\mathbb C)\), hence is simply connected, and its
geometric monodromy is the scalar central element of order three.
It supplies no \(\mathbb Z^4\)- or \(\mathbb Z^6\)-valued first-homology local
system and no order-four peripheral loop.

Second, the trace-one rank-one locus is the homogeneous space
\[
 \mathbb{CP}^2=SU(3)/S(U(1)\times U(2)).
\]
Any \(SU(3)\)-invariant map from it to the candidate base, on which no
\(SU(3)\)-action is part of the construction, is constant.  Choosing a line
\(\ell\in\mathbb{CP}^2\) does define the exact positive compression
\[
 \operatorname{Herm}_3(\mathbb C)_{\ge0}\longrightarrow
 \operatorname{Herm}(\ell^\perp)_{\ge0},
 \qquad H\longmapsto P_{\ell^\perp}H|_{\ell^\perp},     \tag{TEB.53}
\]
but its target is a four-real-dimensional Lorentz cone, not
\(\mathbb P^1\), and it contains no integral torus lattice.  An integral
quotient such as \(P\) requires the additional choice of maximal-torus and
weight-lattice data.  Irreducibility of the fundamental \(SU(3)\)
representation forbids an \(SU(3)\)-equivariant rank-two linear quotient of
\(\mathbb C^3\).

In particular, the homomorphism (TEB.14) belongs to the toroidal degeneration,
not to the Gram cone.  A generic conjugacy orbit in the trace-normalized
Hermitian cone is \(SU(3)/T^2\), which is simply connected, and the determinant
page is simply connected as well.  Neither has first homology
\(\mathbb Z^6\).  Calling any map from it a \(T^6\) filling map before
constructing a torus bundle would be a type error.

\subsection{The remaining global problem}
\label{subsec:teb-remaining}

The local bridge and the finite lattice lift isolate the remaining question.
One must construct a proper complex, or explicitly stratified-real, map
\[
 \widetilde f:\widetilde X^8\longrightarrow\mathbb P^1
\]
and a map over the same base
\[
 \Pi:\widetilde X\longrightarrow X                              \tag{TEB.54}
\]
with all of the following literal properties:
\begin{enumerate}
\item on every regular fibre, \(\Pi\) is the elliptic extension (TEB.31) and
its marked integral map is \(\Phi\);
\item over the cusp disc, \(\Pi\) is the toric morphism (TEB.35), including
the central strata and the filling square (TEB.26);
\item over each finite-order disc, the affine monodromy and translation are
those in (TEB.43)--(TEB.45), and its punctured overlap agrees with the regular
period family;
\item the actual based van Kampen maps and integral Mayer--Vietoris maps are
computed from these inclusions, rather than inferred from the displayed
monodromy matrices.
\end{enumerate}

Equations (TEB.24)--(TEB.45) prove that there is no integral monodromy or
affine-freeness obstruction to these requirements.  They do not supply the
global period and overlap construction in item 3.

There is also an unavoidable topological conclusion.  The rank-three
toroidal central fibre has Euler characteristic six.  The two finite fibres
in (TEB.45) are free finite quotients of \(T^6\), so each has Euler
characteristic zero.  If these are the only three special fibres, Euler
integration gives
\[
 \chi(\widetilde X)=6.                                  \tag{TEB.55}
\]
Therefore this direct complex globalization cannot be \(S^8\), whose Euler
characteristic is two.  Independently, \(S^8\) admits no almost-complex
structure.  A real \(S^8\) construction would have to add non-complex
singular filling data with total Euler contribution \(-4\), then pass the
full fundamental-group, integral-homology, and tangent-clutching tests.

The nine-dimensional Gram cone remains a valid independent
representation-theoretic carrier.  Relating it to (TEB.54) requires a new,
explicit morphism from determinant-rank data to period, lattice, and filling
data.  Neither the determinant, the Cartan cubic, nor the
\(\mathbb{CP}^2\)-family of normal Lorentz cones supplies such a morphism.
